%% titre: Devoir surveillé — Suites et séries
%% niveau: Licence 2
%% matiere: Analyse
%% annee: 2024-2025
%% description: Exemple d'ancien devoir (à remplacer). Durée : 2 h, sans documents.

\begin{exercice}[Suites, 6 points]
Soit $(u_n)_{n\in\N}$ définie par $u_0=1$ et $u_{n+1}=\sqrt{2+u_n}$.
\begin{enumerate}
  \item Montrer que, pour tout $n\in\N$, $1\le u_n\le 2$.
  \item Montrer que $(u_n)$ est croissante.
  \item En déduire que $(u_n)$ converge et déterminer sa limite.
\end{enumerate}
\begin{correction}
\begin{enumerate}
  \item Par récurrence : $u_0=1\in[1,2]$ et si $1\le u_n\le 2$ alors $\sqrt3\le u_{n+1}\le 2$.
  \item $u_{n+1}-u_n=\dfrac{2+u_n-u_n^2}{\sqrt{2+u_n}+u_n}=\dfrac{(2-u_n)(1+u_n)}{\sqrt{2+u_n}+u_n}\ge 0$.
  \item Croissante et majorée, elle converge vers $\ell$ vérifiant $\ell=\sqrt{2+\ell}$, d'où $\ell=2$.
\end{enumerate}
\end{correction}
\end{exercice}

\begin{exercice}[Séries, 7 points]
Étudier la nature des séries de terme général
\[
  a_n=\frac{1}{n^2+n},\qquad b_n=\frac{n!}{n^n},\qquad c_n=\frac{(-1)^n}{\sqrt{n}} .
\]
\begin{correction}
On a $a_n=\frac1n-\frac1{n+1}$ : série télescopique convergente, de somme $1$. Pour $b_n$, $\frac{b_{n+1}}{b_n}=\left(1+\frac1n\right)^{-n}\to e^{-1}<1$ : convergence (d'Alembert). Pour $c_n$, le critère des séries alternées donne la convergence.
\end{correction}
\end{exercice}
